\section{Preliminaries}
\label{sec:preliminaries}

\begin{table}[ht]
\centering
\caption{Notation used throughout the preliminaries.}
\label{tab:notation}
\renewcommand{\arraystretch}{1.18}
\begin{tabularx}{\textwidth}{@{}>{\ttfamily}p{0.27\textwidth} >{\raggedright\arraybackslash}p{0.18\textwidth} X@{}}
\toprule
\normalfont\textbf{Command} & \textbf{Display} & \textbf{Meaning} \\
\midrule
\string\clustering & clustering & A partition of the vertex set \\
\string\clusterings & clusterings & Partitions of the vertex set \\
\string\Clustering & Clustering & Capitalized form of a clustering \\
\string\Clusterings & Clusterings & Capitalized form of clusterings \\
\string\cost & $f$ & Cost function \\
\string\estcost & $\widehat{\cost}$ & Estimated cost \\
\string\card\{X\} & $\card{X}$ & Cardinality of $X$ \\
\texttt{\textbackslash dist\{\textbackslash cls\{C\},\textbackslash cls\{P\}\}} & $\dist{\cls{C},\cls{P}}$ & Distance between the clusterings $\cls{C}$ and $\cls{P}$ \\
\string\weight & $d$ & Edge weight (fraction of input clustering separates $ e $) \\
\string\weight\{e\} & $\weight{e}$ & Weight of edge $e$ \\
\string\estweight & $\widehat d$ & Estimated edge weight \\
\string\estweight\{e\} & $\estweight{e}$ & Estimated weight of edge $e$ \\
\string\mweight & $p$ & Minimum edge weight \\
\string\mweight\{e\} & $\mweight{e}$ & Minimum weight $\min(\weight{e},1-\weight{e})$ of edge $e$ \\
\string\summweight & $B$ & Sum of minimum edge weights $ := \sum_{e}\mweight{e} $\\
\string\summweight\{t\} & $\summweight{t}$ & Sum of minimum edge weights at time $t$ \\
\string\cls\{C\} & $\cls{C}$ & Generic clustering $ \cls{C} $\\
\string\inpcls & $\mathcal P$ & A generic input clustering \\
\string\inpcls\{i\} & $\inpcls{i}$ & $i$-th input clustering \\
\string\outcls & $\mathcal S$ & Output clustering \\
\string\outcls\{t\} & $\outcls{t}$ & Output clustering at time $t$ \\
\string\nupdate & $\nupdate$ & Number of updates \\
\texttt{\textbackslash ncocls\{\textbackslash cls\{C\}\}} & $\ncocls{\cls{C}}$ & Number of co-clustered pairs in $\cls{C}$ \\
\texttt{\textbackslash ncocls\{\textbackslash cls\{C\},\textbackslash cls\{P\}\}} & $\ncocls{\cls{C}, \cls{P}}$ & Number of co-clustered pairs in both $\cls{C}$ and $\cls{P}$ \\
\string\expt & $\mathbb E$ & Expectation operator \\
\string\expt\{X\} & $\expt{X}$ & Expectation of $X$ \\
\bottomrule
\end{tabularx}
\end{table}

\paragraph{\texorpdfstring{\Clusterings}{Clusterings}.}
A \emph{\clustering} of a vertex set $ V $ is a collection of nonempty, pairwise
disjoint subsets, called \emph{clusters}, whose union is that vertex set.
In a \clustering, two vertices are \emph{co-clustered} if they belong to
the same cluster, and \emph{separated} if they belong to different clusters.

For two \clusterings $ \cls{C}, \cls{P} $ of a set $ V $, their distance $ \dist{\cls{C}, \cls{P}}$ is the number of unordered pairs that are co-clustered in exactly one of $ \cls{C},\cls{P} $:
\begin{align}
    \dist{\cls{C}, \cls{P}} = \card{\{ u,v \}\subset V:\ u,\ v \text{ are co-clustered in exactly one of } \cls{C} \text{ or } \cls{P}} \nonumber
\end{align}
For any subset $ V'\subseteq V $, let $ \cls{C}|_{V'} $ denote the restriction of $ \cls{C} $ to $ V' $.

\paragraph{Consensus Clustering Problem.}\label{par:consensus_clustering_} % (fold)
Given $ m $ \clusterings $ \inpcls{1}, \inpcls{2}, \dots, \inpcls{m} $ of a vertex set $ V $ as input, \emph{consensus clustering problem} seeks a \clustering $ \cls{C} $ of $ V $ that minimizes the \emph{cost} function
\begin{align}
    \cost_{ V }(\cls{C}) = \dfrac{1}{m} \sum_{i=1}^{m} \dist{\cls{C}, \inpcls{i}} .\nonumber
\end{align}
The subscript $ V $ is omitted when the vertex set is clear from context.

We denote by $ \OPT = \min_{\cls{C}:\ \cls{C} \text{ is a \clustering\ } \text{ of } V}\cost(\cls{C}) $ the optimal cost of the consensus clustering problem on the given input \clusterings.
% paragraph Consensus Clustering. (end)

\tienlong{Rememeber to use throughout either identifier or labels}
\paragraph{Presenting and Querying a \texorpdfstring{\Clustering}{clustering}.}
Each input \clustering is presented by assigning a cluster label to each vertex.
Two vertices are co-clustered if and only if they have the same label.

\paragraph{Offline Setting.}\label{par:offline} % (fold)
The whole input is given at once.
% paragraph Offline (end)

The algorithm can obtain a vertex's cluster identifier in any specified
input in constant time and can sample a uniform input index in constant
time. Input labels stay fixed while a vertex is active. We treat a
reinserted vertex as a new vertex. We use a word RAM with counters and identifiers stored
in a constant number of words at fixed accuracy. Dynamic perfect
hashing gives worst-case constant lookup time and expected amortized
constant time for insertions and deletions~\cite{Dietzfelbinger1994}.

The offline algorithm outputs one cluster identifier for every vertex.
The online algorithm starts empty, sees vertex insertions, and reports
the new label and all changes to previously published labels. Recourse
counts changes to an individual vertex's published label after it
arrives. The fully dynamic algorithm also starts empty. Its updates
insert or delete a vertex, and it provides an implicit \clustering
through queries. A label query returns a vertex's identifier in the
maintained \clustering. A same-cluster query tests equality of two such
identifiers. When a probability guarantee holds simultaneously over the update
sequence, a single event ensures the approximation or competitive ratio
after every update.

% (end)

In a \emph{non-offline} setting, each vertex is inserted or deleted one at a time.
Let $ \nupdate $ be a positive integer and let $ \{ 1,2,\dots, \nupdate \} $ index each time step where a vertex is inserted or deleted.
For each $ t\in \{ 1,2, \dots, \nupdate \} $, an \emph{active set} $ V_t \subseteq V $ is the set of vertices that have arrived by time $ t $ and have not subsequently been deleted.
We denote by $ \OPT_t $ the optimal cost of the consensus clustering problem on the input \clusterings restricted to $ V_t $.
A clustering $ \cls{C} $ of $ V_t $ is $ \alpha $-approximation to $ \OPT_t $ if $ \cost_{V_{t}}(\cls{C}) \leq \alpha \OPT_{t} $.
An \emph{expected approximation ratio} $\alpha$ means $\expt{\cost_{t}(C_t)}\le\alpha\OPT_t$. \tienlong{Specify later what is the randomness in the expectation.}

In this paper, we consider two \emph{non-offline} settings: 

\paragraph{Online with Recourse.}\label{par:online_} % (fold)
Each vertex arrives one at a time, revealing upon arrival all its labels in the input \clusterings.

\tienlong{Recourse is defined as the number of times a vertex's label changes between two consecutive time steps.}
% paragraph Online. (end)

\paragraph{Fully Dynamic}\label{par:fully_dynamic} % (fold)
Nodes can be inserted or deleted one at a time. Each vertex reveals upon arrival all its labels in the input \clusterings.

\tienlong{Update time is the time the algorithm takes to maintain the solution}
\tienlong{Query time is the time the algorithm takes to answer a query}
% paragraph Fully Dynamic (end)

For a pair $e=\{u,v\}$, let $ \weight{e}= \weight{uv} $ be the fraction of inputs
that separate its endpoints. The pair contributes $ \weight{e} $ to $ \cost(\cls{C}) $ if
$ \cls{C} $ assigns its endpoints to the same cluster, and $ 1- \weight{e} $ otherwise. Its minimum possible contribution
is $ \mweight{e}= \mweight{uv}=\min\{\weight{e},1-\weight{e}\}$. The sum $ \summweight=\sum_e \mweight{e} $, over unordered
pairs, is therefore a lower bound on $\OPT$. The values $ \weight{uv} $ satisfy
the triangle inequality because each input's separation indicators do.

On insertion-only streams, $\summweight{t}$ and $\OPT_t$ are nondecreasing. The
majority decision for a pair assigns its endpoints to the same cluster
if $\weight{e}<1/2$ and to different clusters if $\weight{e}>1/2$. If $\summweight{t}<1/2$, there are no ties. The expected number of
pairs on which a uniform input disagrees with the majority decisions
is then $\summweight{t}<1$, so some input realizes all of them. That \clustering
attains $\OPT_t$. Label each cluster by its smallest arrival index to
preserve all old labels while $B_t<1/2$.


\paragraph{Maintaining \clustering distances.}
For \clusterings $C,P$, let $ \ncocls{C}$ count the pairs co-clustered in $C$
and let $ \ncocls{C,P}$ count the pairs co-clustered in both $C$ and $P$. Then
$\dist{C,P}=\ncocls{C}+\ncocls{P}-2\ncocls{C,P}$. Each term is the sum of $\binom b2$
over the sizes $b$ of the clusters in $C$, the clusters in $P$, or
their intersections, respectively. We can therefore evaluate one
distance in $O(n)$ expected time or update it in expected amortized
$O(1)$ time per insertion, deletion, or change of one vertex's label.
We use the triangle inequality in the forms
$|\dist(C,P)-\dist{A,P}|\le \dist{C,A}$ and $\dist{C,A}\le \cost(C)+\cost(A)$.

\subsection{Sampling input \clusterings and selecting a candidate}
\label{sec:ensemble-transfer}

Unless specified otherwise, sample input indices independently and
uniformly with replacement. Repetitions are allowed, and we retain the
order of the sampled indices. Associate each vertex with the ordered tuple of its cluster identifiers
in the corresponding input \clusterings. Let
$\estcost$ denote the average cost over these \clusterings, and let
$\estweight{e}$ be the fraction of them separating pair $e$. A
\emph{candidate} is a \clustering produced by a run of an algorithm.
To compare candidates, independently sample further input \clusterings
and estimate the candidates' average distances to them. Only the pairs
on which two candidates disagree can affect the error in their
estimated cost difference.

\begin{lemma}\label{lem:sample-difference}
On a fixed vertex set $V$, let $\estweight{e}$ denote the separation
fraction in a sample of $k$ independently sampled inputs. For $\delta>0$,
let
\[
 R(V)=\sum_{e\subseteq V,\ |e|=2}(2|\estweight{e}-\weight{e}|-\delta)_+,
 \qquad (x)_+=\max\{x,0\}.
\]
For any subset $W\subseteq V$ and \clusterings $C,F$ of $W$,
\[
 \big|(f_W(C)-f_W(F))-(\estcost_W(C)-\estcost_W(F))\big|
 \le \delta \dist{C,F}+R(V).
\]
Moreover, $\expt{R(V)}\le\operatorname{OPT}_V/(k\delta)$.
\end{lemma}
\begin{proof}
Only pairs on which $C$ and $F$ disagree contribute to the difference.
Each contributes at most $2|\estweight{e}-\weight{e}|$ in absolute value.
Apply $(2|x|-\delta)_+\le x^2/\delta$ and the variance of each sampled
fraction to obtain
\[
 \expt{R(V)}\le\frac{1}{k\delta}\sum_e \weight{e}(1-\weight{e})
 \le\frac{\operatorname{OPT}_V}{k\delta}.
\]
The last inequality uses $\weight{e}(1-\weight{e})\le \mweight{e}$ for each pair.
\end{proof}

The next lemma bounds the expected approximation ratio on all input
\clusterings in terms of the ratio on a sample.

\begin{lemma}\label{thm:ensemble-transfer}
Let $1\le\gamma\le2$ and $0<\tau\le1$. If an algorithm has conditional
expected approximation ratio $\gamma$ on every collection of
$k\ge25/\tau^2$ sampled input \clusterings, its expected approximation ratio on the original instance
is at most $\gamma+\tau$.
\end{lemma}
\begin{proof}
Fix a fixed optimal \clustering $F$. Lemma~\ref{lem:sample-difference} bounds the
error even when $C$ depends on the sampled input \clusterings. The \clustering triangle
inequality gives $d(C,F)\le\estcost(C)+\estcost(F)$.
Taking expectations and using $\expt{\estcost(F)}=\operatorname{OPT}$
yields
\[
 \expt{\cost(C)}\le
 \bigl(\gamma+\delta(\gamma+1)+1/(k\delta)\bigr)\operatorname{OPT}.
\]
Take $\delta=\tau/6$ and use the bound on $k$.
\end{proof}

We will also use a deterministic family of short sequences of input indices. Pairwise
independence suffices in Lemma~\ref{thm:ensemble-transfer}, since its
proof uses only the first two moments of the separation fractions.

\begin{lemma}\label{lem:finite-lists}
For any positive integer $k$, one can enumerate $(m k!)^2$ sequences of
$k$ input indices such that a uniformly chosen sequence has pairwise
independent, uniformly distributed entries.
\end{lemma}
\begin{proof}
Set $w=k!$ and $N=mw$. Enumerate $a,b$ modulo $N$ and use input
$1+\lfloor((a+ib)\bmod N)/w\rfloor$ at position $0\le i<k$.
Conditional on $a+ib$ modulo $N$, the residue $a+jb$ is uniform
in a congruence class modulo $\gcd(|j-i|,N)$. If $i\ne j$, this
divisor divides $w$, so each interval of $w$ residues contains the
same number of elements of that class.
\end{proof}

Generate candidates using independently sampled input \clusterings and
independent rounding. To compare them, independently sample further
input \clusterings and use the same \clusterings to estimate all candidate
costs. The resulting cost-difference estimates suffice to select the
candidate with the smallest estimated cost.

\begin{lemma}\label{lem:select-simple}
Suppose each trial has expected approximation ratio $\gamma$ at
each of $U\ge1$ fixed instances, where $1\le\gamma\le2$.
Fix $0<\delta\le1/12$ and $0<\zeta<1$. Run
$T=\lceil (3/\delta)\log(2U/\zeta)\rceil$ independent trials and independently
sample $R=\lceil2\delta^{-2}\log(4UT^2/\zeta)\rceil$ input \clusterings
to estimate costs for all trials and instances. Selecting the candidate of smallest estimated cost gives
approximation ratio $\gamma+6\delta$ at all instances simultaneously
with probability at least $1-\zeta$.
\end{lemma}
\begin{proof}
At an instance with $\OPT>0$, Markov's inequality gives a trial
probability at least $\delta/(\gamma+\delta)\ge\delta/3$ of producing
a candidate of cost at most $(\gamma+\delta)\OPT$. With probability
at least $1-\zeta/2$, some candidate $C_\star$ meets this bound at each
instance. When $\OPT=0$, every candidate has cost zero almost surely.

Condition on all candidate \clusterings. The cost difference between two candidates
$C,A$ on one input lies in $[-\dist{C,A},\dist{C,A}]$. By Hoeffding's inequality
and a union bound over at most $UT^2$ comparisons, every average cost
difference has error at most $\delta \dist{C,A}$ except with probability
$\zeta/2$. In particular, for the candidate $C$ with smallest estimated
cost, $\cost(C)-\cost(C_\star)\le\delta \dist{C,C_\star}\le\delta(\cost(C)+\cost(C_\star))$.
The resulting factor is at most
$(1+\delta)(\gamma+\delta)/(1-\delta)\le\gamma+6\delta$.
\end{proof}

The preceding rule is sufficient with constant success probability. The next
lemma gives the full dependence on the offline failure probability using a linear number
of comparisons in $\log(2/\zeta)$, without comparing every pair of
candidates in a large sample.

\begin{lemma}\label{lem:offline-confidence}
Let $1\le\gamma\le2$ and $0<\varepsilon\le1$. An offline algorithm
with approximation ratio $\gamma+\varepsilon/2$ and success probability
$9/10$ can be amplified to approximation ratio $\gamma+\varepsilon$
and success probability $1-\zeta$ using
$O(\log(2/\zeta))$ expected independent calls to that algorithm and
$O(\varepsilon^{-2}\log(2/\zeta))$ expected \clustering-distance
evaluations.
\end{lemma}
\begin{proof}
In each attempt, draw a candidate $C$ and $M=\Theta(\log(2/\zeta))$
independent \clusterings from the given algorithm. For each resulting \clustering $Q$, let
$D=\dist{C,Q}$. If $D>0$, average
$(\dist{C,P}-\dist{Q,P})/D$ over $O(\varepsilon^{-2})$ independently sampled
uniform input \clusterings $P$. The comparison accepts the candidate
if this average is at most $\theta=\varepsilon/28$ or if $D=0$.
Accept $C$ if at least half of its comparisons succeed. Otherwise
start an independent attempt.

Each normalized difference lies in $[-1,1]$. When $\cost(C)\le \cost(Q)$,
one comparison accepts the candidate with probability at least $19/20$, by
Hoeffding's inequality. If $\cost(C)>(\gamma+\varepsilon)\OPT$ and
$\cost(Q)\le(\gamma+\varepsilon/2)\OPT$, the mean difference is at least
$\varepsilon/(4\gamma+3\varepsilon)>2\theta$, since
$D\le \cost(C)+\cost(Q)$. The mean is at least one when $\OPT=0$ in this
case. Thus a candidate above the target is accepted in one comparison
with probability at most $1/10+(9/10)/20=29/200$.

With probability at least $1/4$, $\cost(C)$ is at or below the lower
quartile of the output cost distribution. Conditional on such a
candidate, a comparison accepts the candidate with probability at least
$(3/4)(19/20)=57/80$. Choose the constant in $M$ so that
Hoeffding's inequality gives unconditional acceptance probability
at least $1/8$ and probability at most $\zeta/8$ of accepting a
candidate above the target in any one attempt. Independent attempts
then take at most eight repetitions in expectation and have eventual
error at most $\zeta$. Counting the calls and evaluations gives the
bounds.
\end{proof}
